% !TEX program = xelatex
% The Frozen Row — a plain-English picture book about how a terminal
% AI assistant never scrambles your screen history.
% Build:  xelatex -interaction=nonstopmode frozen-row.tex  (run twice for TOC)
\documentclass[11pt,twoside]{book}

% ---------- Fonts (XeLaTeX + fontspec) ----------
\usepackage{fontspec}
\newcommand{\tgpath}{/usr/share/texmf-dist/fonts/opentype/public/tex-gyre/}
\setmainfont{texgyrepagella-regular.otf}[
  Path=\tgpath,
  BoldFont=texgyrepagella-bold.otf,
  ItalicFont=texgyrepagella-italic.otf,
  BoldItalicFont=texgyrepagella-bolditalic.otf,
  Numbers=OldStyle,
  Ligatures=TeX ]
\setsansfont{texgyreheros-regular.otf}[
  Path=\tgpath,
  BoldFont=texgyreheros-bold.otf,
  ItalicFont=texgyreheros-italic.otf,
  BoldItalicFont=texgyreheros-bolditalic.otf,
  Scale=0.92 ]
\newfontfamily\codefont{JetBrainsMono Nerd Font}[
  Scale=0.82,
  Ligatures=NoCommon ]

% ---------- Page geometry ----------
\usepackage[
  paperwidth=6.5in, paperheight=9.25in,
  inner=0.85in, outer=0.7in, top=0.85in, bottom=0.85in,
  headsep=14pt, footskip=26pt ]{geometry}

% ---------- Colors ----------
\usepackage[table]{xcolor}
\definecolor{ink}{HTML}{1A1A22}
\definecolor{accent}{HTML}{2C4A78}      % deep slate blue
\definecolor{accent2}{HTML}{9B4A2F}     % rust
\definecolor{codebg}{HTML}{F5F5F1}
\definecolor{coderule}{HTML}{D9D9CE}
\definecolor{noteblue}{HTML}{EAF0F7}
\definecolor{notebluerule}{HTML}{2C4A78}
\definecolor{warmbg}{HTML}{FBF1E9}
\definecolor{warmrule}{HTML}{9B4A2F}
\definecolor{greenbg}{HTML}{ECF2E8}
\definecolor{greenrule}{HTML}{4A6B34}
\definecolor{rulegray}{HTML}{B8B8B0}
\color{ink}

% ---------- Microtype, spacing ----------
\usepackage[protrusion=true]{microtype}
\usepackage{setspace}
\setstretch{1.08}
\usepackage{ragged2e}
\usepackage[hidelinks]{hyperref}
\hypersetup{
  pdftitle={The Frozen Row: how a terminal keeps your screen from scrambling},
  pdfauthor={A plain-English guide},
  colorlinks=true, linkcolor=accent, urlcolor=accent2 }

% ---------- Headings (titlesec) ----------
\usepackage{titlesec}
\titleformat{\chapter}[display]
  {\sffamily\bfseries\color{accent}}
  {\fontsize{13}{13}\selectfont\MakeUppercase{Chapter \thechapter}}
  {10pt}
  {\fontsize{23}{27}\selectfont\color{ink}}
  [\vspace{4pt}{\color{rulegray}\titlerule[1pt]}]
\titlespacing*{\chapter}{0pt}{6pt}{26pt}

\titleformat{\section}
  {\sffamily\bfseries\large\color{accent}}{}{0pt}{}
\titlespacing*{\section}{0pt}{14pt}{5pt}

% ---------- Running heads / feet ----------
\usepackage{fancyhdr}
\pagestyle{fancy}
\fancyhf{}
\renewcommand{\headrulewidth}{0pt}
\renewcommand{\footrulewidth}{0pt}
\fancyhead[LE]{\sffamily\footnotesize\color{rulegray}\thepage\quad\textbar\quad The Frozen Row}
\fancyhead[RO]{\sffamily\footnotesize\color{rulegray}\leftmark\quad\textbar\quad\thepage}
\renewcommand{\chaptermark}[1]{\markboth{#1}{}}
\fancypagestyle{plain}{\fancyhf{}\fancyfoot[C]{\sffamily\footnotesize\color{rulegray}\thepage}}

% ---------- Lists ----------
\usepackage{enumitem}
\setlist[itemize]{leftmargin=1.4em,itemsep=3pt,topsep=4pt}
\setlist[description]{leftmargin=0pt,style=nextline,font=\sffamily\bfseries\color{accent}}

\usepackage[most]{tcolorbox}

% callout: the big idea (blue)
\newtcolorbox{insight}[1][]{
  colback=noteblue, colframe=notebluerule,
  boxrule=0pt, leftrule=2.5pt, arc=1pt,
  left=10pt, right=10pt, top=7pt, bottom=7pt,
  before skip=10pt, after skip=10pt,
  fonttitle=\sffamily\bfseries\color{notebluerule},
  title={#1},
}
% callout: watch out (warm)
\newtcolorbox{trap}[1][]{
  colback=warmbg, colframe=warmrule,
  boxrule=0pt, leftrule=2.5pt, arc=1pt,
  left=10pt, right=10pt, top=7pt, bottom=7pt,
  before skip=10pt, after skip=10pt,
  fonttitle=\sffamily\bfseries\color{warmrule},
  title={#1},
}
% callout: try it yourself / picture this (green)
\newtcolorbox{pic}[1][]{
  colback=greenbg, colframe=greenrule,
  boxrule=0pt, leftrule=2.5pt, arc=1pt,
  left=10pt, right=10pt, top=7pt, bottom=7pt,
  before skip=10pt, after skip=10pt,
  fonttitle=\sffamily\bfseries\color{greenrule},
  title={#1},
}

% drop cap
\usepackage{lettrine}
\renewcommand{\LettrineFontHook}{\color{accent}\sffamily\bfseries}
\setcounter{DefaultLines}{3}

% ---------- TOC styling ----------
\usepackage{titletoc}
\usepackage{tocloft}
\renewcommand{\cfttoctitlefont}{\sffamily\bfseries\Large\color{accent}}
\renewcommand{\cftchapfont}{\sffamily\bfseries\color{ink}}
\renewcommand{\cftchappagefont}{\sffamily\color{ink}}
\renewcommand{\cftchapleader}{\color{rulegray}\cftdotfill{\cftdotsep}}
\setlength{\cftbeforechapskip}{4pt}

\usepackage{graphicx}

% ================================================================
\begin{document}
\frenchspacing
\emergencystretch=3em
\hbadness=4000
\tolerance=2000

% ---------------- TITLE PAGE ----------------
\thispagestyle{empty}
\begin{titlepage}
\raggedright
\vspace*{1.1in}
{\color{rulegray}\rule{\linewidth}{1pt}}\\[10pt]
{\sffamily\bfseries\fontsize{40}{44}\selectfont\color{ink} The Frozen Row}\\[14pt]
{\sffamily\fontsize{15}{20}\selectfont\color{accent}
  How a chat program keeps your screen\\ from turning into a scrambled mess}\\[10pt]
{\color{rulegray}\rule{\linewidth}{1pt}}\\[24pt]
{\itshape\fontsize{12}{16}\selectfont
A picture-book explanation, in plain English,\\
of one clever little idea \\
and the careful rules that make it work.}\\[10pt]
\vfill
{\sffamily\small\color{accent}
No computer knowledge needed.\\
If you have ever used a receipt printer or a whiteboard,\\
you already know enough to understand this whole book.}\\[20pt]
{\color{rulegray}\rule{\linewidth}{0.5pt}}\\[6pt]
{\small A friendly guide to the trickiest part of a terminal AI assistant.}
\end{titlepage}

% ---------------- FRONT MATTER ----------------
\frontmatter
\pagestyle{plain}

\chapter*{Before we start}
\addcontentsline{toc}{chapter}{Before we start}
\markboth{Before we start}{}

\lettrine{T}{his} is a book about a problem that sounds tiny and turns out
to be surprisingly deep. A computer program is drawing words on a screen,
very fast, and it must never make a mistake. That's it. That's the whole
problem.

But there is a catch that makes it hard, and the catch is the heart of the
whole story. So before anything else, here is a promise: \textbf{you will
not need to know anything about computers to understand this book.} Every
idea in here is explained with things you already know---a receipt
printer, a whiteboard, a notebook, a careful painter, a librarian. Whenever
something has a fancy computer name, we will give it a nickname and use the
nickname instead.

The program we are talking about is a kind of chatbot that lives in a
plain text window on a computer---the sort programmers use. You type a
question, and an AI types back an answer, live, one word at a time, like a
person texting you. The answer can get long. Long enough to scroll off the
top of the screen. And later you can scroll back up to re-read it.

Simple, right? You'd think so. But making that work \emph{without ever
scrambling the screen} took years of careful thought and a small pile of
clever rules. This book explains those rules, one at a time, from the very
beginning.

Read it front to back like a story. Each chapter builds on the one before.
By the end you'll understand something most professional programmers never
stop to think about---and you'll see why it's genuinely beautiful.

\begin{pic}[How to read the coloured boxes]
Throughout the book you'll see three kinds of coloured boxes.
\textcolor{notebluerule}{\textbf{Blue boxes}} hold the big idea of a
chapter---the one sentence to remember.
\textcolor{warmrule}{\textbf{Orange boxes}} are warnings: a tempting
mistake and why it goes wrong.
\textcolor{greenrule}{\textbf{Green boxes}} ask you to picture something,
so the idea sticks.
\end{pic}

\tableofcontents

% ---------------- MAIN MATTER ----------------
\mainmatter
\pagestyle{fancy}

% ======================================================
\chapter{The receipt printer}
% ======================================================
\lettrine{I}{magine} a receipt printer---the kind at a shop checkout. Paper
feeds up from a slot, and as it prints, the paper rolls upward and out the
top. New lines appear at the bottom; old lines climb toward the top and
eventually disappear over the edge.

Now here is the important part. Once a line of paper has rolled up and out
of the machine, \textbf{you cannot un-print it.} It's done. If there was a
typo on that line, too bad---it has left the building. You could tear off
the whole roll, but you cannot reach back in and fix that one line. It is,
in a word, \emph{frozen}.

A computer screen in a text window works almost exactly like this receipt
printer. As new text appears at the bottom, old text scrolls up and off the
top. And---this is the whole book in one sentence---\textbf{once a line has
scrolled off the top of the screen, the program can never change it again.}

You, the human, can still \emph{see} those old lines: you scroll up with
your mouse wheel and there they are, saved in a big roll of history. But the
program that put them there has lost all power over them. They are frozen,
exactly like a printed receipt.

\begin{pic}[Picture this]
Hold a paper receipt in your hand. Try to change the third line from the
top without tearing anything or scribbling on it. You can't. That
helplessness---that's what the program feels about every line that has
scrolled off the screen.
\end{pic}

So why does this matter so much? Because our chat program is writing to the
screen \emph{constantly}. Every time a new word of the AI's answer arrives,
the program redraws things, and a little more text scrolls up and freezes.

That means the program gets exactly \textbf{one chance} to get each line
right: the split-second before that line scrolls off the top. If it puts the
wrong thing there in that instant, the mistake is frozen into your history
forever, and you'll see it every time you scroll up.

\begin{insight}[The one big idea]
Every line eventually scrolls off the top and freezes. The program has one
chance to get each line right---the moment before it freezes. Everything
else in this book exists to make sure that, at that exact moment, the line
is correct.
\end{insight}

That is why the book is called \emph{The Frozen Row}. A ``row'' is just one
line of text on the screen. The entire design is about one thing: making
sure that when a row freezes, it's right.

% ======================================================
\chapter{Why not just take over the whole screen?}
% ======================================================
\lettrine{H}{ere's} a fair question. Some programs \emph{don't} have this
problem at all. Think of a full-screen program---a text editor, or one of
those live system monitors with the moving bars. They take over the
\emph{entire} screen. Nothing scrolls off; there's no receipt roll; when you
quit, the whole thing vanishes and your old screen pops back exactly as it
was.

If that's so much easier, why doesn't our chat program do the same thing?

Because it would ruin the thing people love most about it. Our chat program
deliberately shares the screen with everything else you were doing. Your
earlier commands, the AI's answers, more commands, more answers---all of it
piles up together in one long roll of history that you can scroll through
later, like a diary of your whole session. It feels natural. It feels like
part of your terminal, not a separate app that hijacks the screen.

The polite word for this is \textbf{inline}: the program draws its stuff
\emph{in line} with everything else, right where your cursor already is,
instead of taking over.

\begin{pic}[Two kinds of program]
A full-screen program is like a \textbf{whiteboard}: it owns the whole
surface, wipes it clean when it starts, and you can erase and redraw any
part any time. Our inline program is like \textbf{writing at the bottom of a
shared notebook} that already has other people's notes above yours---and the
pages, once turned, are glued shut.
\end{pic}

So the chat program chose the harder life on purpose. It gets the lovely
shared-history feeling, but it pays for it: it has to live with the receipt
printer, and the frozen rows, and the one-chance-to-get-it-right rule. The
rest of this book is how it pays that price without ever slipping up.

When it \emph{does} slip up---and early versions did---you get ugly things:
the little text box you type into suddenly appears twice, one copy stranded
higher up. Or a whole answer gets accidentally printed twice into your
history. Or your older commands slowly creep upward for no reason. Every one
of those bugs is the same underlying mistake: the program got confused about
what was really on the screen, and a wrong row froze.

% ======================================================
\chapter{The magic notebook}
% ======================================================
\lettrine{N}{ow} we reach the single cleverest idea in the whole design, and
happily it's an easy one to picture.

Remember that the program can never look at the screen. This sounds
unbelievable, but it's true: the program can \emph{send} text to the screen,
but it can never \emph{read back} what's currently showing. It's like
mailing letters into a house with no windows---you can send things in, but
you can never peek to see what's inside.

So how can the program possibly know what to change, if it can't see what's
there? Here's the trick. The program keeps a \textbf{second copy}---a
private notebook---in which it writes down \emph{exactly what it believes is
on the screen right now}. Every time it sends something to the real screen,
it also updates its notebook to match. The notebook is the program's memory
of the screen.

We'll call this notebook the \textbf{shadow}, because it's meant to be a
perfect shadow of the real screen---moving in step with it, always
identical.

\begin{insight}[The shadow]
The program can't see the screen, so it keeps a private notebook---the
shadow---that says what it thinks is on the screen. The single most
important rule in the entire design is: \textbf{the shadow must always
match the real screen, exactly.}
\end{insight}

Why does this matter? Because of how the program decides what to redraw. When
a new word of the answer arrives, the program figures out the new picture it
wants, then \textbf{compares} it against the shadow---its notebook of what's
already there. Wherever they differ, it sends just those changes to the
screen. Wherever they're the same, it sends nothing (why repaint a line
that's already correct?).

This comparing step is the beating heart of the program, and we'll call it
the \textbf{compare}. It's efficient and lovely: even if the answer is ten
thousand lines long, only the few lines that actually changed get touched.

But notice: the compare trusts the notebook completely. If the notebook is
\emph{wrong}---if it says a line holds one thing but the real screen shows
another---then the compare makes bad decisions, sends text to the wrong
places, and sooner or later a wrong line freezes. Everything downstream
depends on the shadow being a perfect match.

\begin{pic}[The two ways it can go wrong]
Suppose the notebook is wrong about a line. Two bad things can happen. Either
the program thinks ``this line is already fine'' when it isn't---so it leaves
a stale, wrong line on screen (things look \emph{stuck}). Or it thinks ``this
line changed'' when it didn't---so it needlessly repaints, and every needless
repaint is a chance to nudge something out of place. Both roads end at a
frozen wrong row.
\end{pic}

Almost every bug this program ever had, in its whole history, was one thing:
the shadow briefly stopped matching the real screen. Keep that in mind. It's
the villain of the entire story, and every rule from here on is armour
against it.

% ======================================================
\chapter{Keeping the notebook honest}
% ======================================================
\lettrine{S}{o} the shadow---the notebook---must always match the real
screen. How does the program keep that promise? With a simple, disciplined
habit, repeated on every single update.

The habit is this: \textbf{whenever the program paints something onto the
real screen, in the very same breath it writes that same thing into the
notebook.} Never one without the other. Paint a line; copy that line into
the notebook. Paint another; copy that too. The two actions are welded
together, like a cashier who can't hand you an item without also scanning it.

If it follows this habit perfectly, then by simple logic the notebook can
never fall out of step. Think about it: the lines it didn't touch this time
were already correct (they matched last time, and nothing changed them). The
lines it \emph{did} touch, it just copied into the notebook as it painted
them. So after every update, top to bottom, notebook equals screen. Always.

\begin{pic}[Painter and photograph]
Imagine a careful painter touching up a mural. Every time she paints a
stroke on the wall, she immediately makes the identical stroke on a small
photograph of the wall she keeps in her pocket. She never looks at the wall
directly---she's not allowed---but because she updates the photo with every
stroke, her photo is always a perfect copy of the wall. The notebook is that
pocket photograph.
\end{pic}

There's one more thing the program has to be careful about: knowing where
its ``pen'' is. On the real screen, the program moves an invisible pen
(programmers call it the cursor) and text comes out wherever the pen is. But
remember---it can't look at the screen, so it can't see where the pen
actually is. It has to \emph{remember} where it left the pen last time, and
trust that memory.

This is where a lot of the danger lives. If the program's memory of ``where
the pen is'' drifts even one line away from reality, then the next thing it
draws lands one line off, and that off-by-one mistake can freeze. So the
program is extremely disciplined about the pen: every move is measured
carefully from where it believes the pen already is, and every action that
could move the pen is accounted for.

\begin{insight}[The habit that keeps the promise]
Paint and record in the same motion. Track the pen with obsessive care. Do
those two things faithfully and the notebook stays a perfect shadow of the
screen---which is the one promise everything else depends on.
\end{insight}

% ======================================================
\chapter{One line at a time}
% ======================================================
\lettrine{L}{et's} slow down and actually watch the program do one update, in
plain steps. A new word has arrived; the program wants to update the screen.
Here's what it does, told as a little recipe.

\textbf{Step one: figure out what should be on screen now.} It lays out the
new picture---all the text, wrapped and arranged---in a fresh workspace. Call
this the \emph{new picture}.

\textbf{Step two: find the first line that's different.} It goes down the
new picture and the notebook side by side, top to bottom, looking for the
first line where they disagree. Everything above that first difference is
already correct and can be left completely alone. On a normal update, the
only thing that changed is a bit more text at the very bottom---so this first
difference is way down near the end, and almost the whole screen gets
skipped. That's why it's fast.

\textbf{Step three: move the pen there.} It moves its invisible pen up or
down from where it remembers leaving it, to reach that first changed line.
This is the careful pen-tracking from the last chapter.

\textbf{Step four: repaint from there to the bottom.} Going line by line, it
paints each changed line, and---following the sacred habit---copies each one
into the notebook as it goes. Lines that happen to be unchanged in the
middle, it just steps past without repainting.

\textbf{Step five: if the new picture is shorter, tidy up.} Sometimes the new
picture is \emph{shorter} than before (say the little type-box shrank). Then
there are leftover lines below the new content still showing old text, so the
program erases them.

And that's a full update. New word in, a few lines repainted, notebook kept
in sync, done---often in well under a thousandth of a second, dozens of times
a second, for as long as the answer keeps streaming in.

\begin{pic}[The whole loop in one image]
A word arrives. The painter glances at her pocket photo, spots the one place
the wall should change, walks her brush there, paints the change, dabs the
same change onto the photo, and steps back. Over and over, faster than you
can watch. The wall and the photo move as one.
\end{pic}

There's a subtle bit in step four worth flagging, because it's where the
frozen-row danger becomes real. To grow the picture and make room for new
text at the bottom, the program pushes everything up---and pushing up is
exactly what makes the top line scroll off and \emph{freeze}. So the moment
of growth is the moment of freezing. That's the dangerous instant, and the
next chapter is entirely about defending it.

% ======================================================
\chapter{The last-second double-check}
% ======================================================
\lettrine{H}{ere} is the scariest instant in the whole program: the moment a
line is about to scroll off the top and freeze forever. If the notebook was
even slightly wrong about that line---maybe a rare hiccup left it stale---then
the wrong text is about to become permanent. And unlike every other mistake,
this one can never be fixed afterward.

The program's response is beautifully paranoid. For any line that is
\emph{about to freeze this instant}, it doesn't trust the compare at all. It
just \textbf{repaints that whole line from scratch}, start to finish, even if
the notebook swears the line is already fine.

Think of it as a bouncer at the exit. Most lines leave the screen quietly.
But for the handful of lines stepping out the top door \emph{right now}---the
ones about to freeze---the bouncer stops each one and says, ``Let me redraw
you completely, just to be certain, because after this door there's no coming
back.'' It costs a few extra keystrokes of work. It's worth it, because the
alternative is a permanent scar in your history.

\begin{insight}[Belt and braces]
The notebook \emph{should} always be correct---but freezing is forever, so
the program refuses to gamble on the lines about to freeze. It repaints them
in full at the last second, guaranteeing they're right no matter what. A
little wasted effort buys a promise that can never be broken.
\end{insight}

Why not do this for \emph{every} line, then, and skip the whole compare? Two
reasons. First, it would be wasteful---repainting the entire screen on every
tiny change, thousands of times, would be slow and make the text flicker.
Second, and more importantly, most lines aren't in danger; they're safely in
the middle of the screen with no risk of freezing this instant. The genius is
in aiming the expensive, careful repaint at \emph{exactly} the lines at
risk---the ones at the top edge about to scroll away---and trusting the fast
compare everywhere else.

\begin{pic}[Where the danger lives]
Picture the screen as a room with a trapdoor at the top. Lines near the
bottom are safe. Lines drifting toward the trapdoor are in danger. Lines
actually stepping onto the trapdoor this second are the ones the program
double-checks by hand. Everyone else, it waves through.
\end{pic}

% ======================================================
\chapter{A quick way to check the notebook}
% ======================================================
\lettrine{T}{here's} a nagging worry hiding under everything so far. We keep
saying ``as long as the notebook is correct.'' But what if, somehow, it
\emph{isn't}? What if some rare glitch scribbled in the notebook when nobody
was looking? Before the program trusts the notebook for another update, it
would be nice to quickly confirm the notebook hasn't been corrupted.

The obvious way is to re-read the whole notebook, line by line, and check
it's all still sensible. But the notebook can be thousands of lines long, and
doing that on \emph{every} update, dozens of times a second, would be
painfully slow.

So the program uses a shortcut you might recognise from everyday life: a
\textbf{tally}. Alongside the notebook, it keeps a single little number---call
it the \emph{fingerprint}---computed from all the lines. If even one letter in
the notebook changes, the fingerprint changes too. So instead of re-reading
thousands of lines, the program just glances at the fingerprint. Matches what
it should be? The notebook is intact. Doesn't match? Something's wrong---stop,
don't trust it.

\begin{pic}[The jar of marbles]
Imagine a jar of marbles and a sticky note on the lid saying how many are
inside. To check nobody's stolen one, you don't have to count the marbles---you
just read the note and shake the jar. The fingerprint is that sticky note: a
tiny summary that instantly reveals if anything's off.
\end{pic}

And there's an extra bit of cleverness. When the program repaints a few lines,
it doesn't recompute the fingerprint from scratch (that would mean re-reading
everything again). Instead it \emph{adjusts} the tally: it subtracts out the
old lines' contribution and adds in the new lines'. Like updating a running
total on a receipt---you don't re-add every item, you just add the newest one.
So keeping the fingerprint current costs almost nothing, no matter how long
the answer gets.

\begin{insight}[Cheap, constant vigilance]
A single fingerprint number lets the program verify its notebook is intact in
an instant, on every update, without ever re-reading the whole thing. Cheap
enough to do always---so it always does.
\end{insight}

If that fingerprint ever fails to match, the program knows its notebook can no
longer be trusted, and it switches into a special repair mode. We'll meet those
repair modes later. For now, the point is simply: the program is constantly,
cheaply checking its own memory---and it never draws based on a notebook it
hasn't just confirmed.

% ======================================================
\chapter{Rules the program cannot break}
% ======================================================
\lettrine{U}{p} to now, every safeguard has been a \emph{habit}---a discipline
the program follows. Paint and record together. Double-check the freezing
lines. Verify the fingerprint. Good habits. But here's an uncomfortable truth
the people who built this program learned the hard way:

\begin{quote}
\itshape Almost every bug in the program's whole history was a moment when
someone, somewhere, \emph{forgot} to follow one of the habits.
\end{quote}

Habits can be forgotten. A new piece of code takes a shortcut, skips the
fingerprint check ``just this once,'' and months later a user sees a scrambled
screen. So the builders did something clever: they stopped relying on
discipline, and made the computer itself \emph{refuse} to let the habits be
skipped.

Here's the idea, and it's genuinely elegant. Think of a nightclub that will
not let you onto the dance floor without a wristband, and the \emph{only} way
to get a wristband is to pass through the security check at the door. There's
no back entrance. You literally cannot be on the dance floor without having
been checked---not because everyone's being careful, but because the building
is shaped so that skipping the check is impossible.

The program does exactly this. The step that draws to the screen demands a
kind of ``wristband''---a token that proves the fingerprint was just checked
and the notebook is confirmed good. And the only way to get that token is to
actually run the check. So ``draw without checking the notebook first'' isn't
just discouraged---it's \emph{impossible}. The program won't even assemble if
you try; it's rejected before it ever runs.

\begin{insight}[From good habits to unbreakable rules]
Discipline can be forgotten. So the design makes the safety checks impossible
to skip: you cannot reach the draw step without first passing the check that
hands you the permission token. The bug of ``forgetting to check'' simply
cannot exist anymore.
\end{insight}

These tokens come as a matched set of wristbands, and the next two short
chapters describe them. But the pattern is always the same: to do the
dangerous thing, you must first hold the proof that it's safe---and the proof
can only be earned honestly.

% ======================================================
\chapter{The two wristbands}
% ======================================================
\lettrine{T}{here} are two things that must be true before the program is
allowed to draw, and so there are two wristbands---two proofs it must hold.

\textbf{Wristband one: ``the notebook matches the visible screen.''} This is
the fingerprint check from Chapter 7, turned into a token. You run the check;
if it passes, you receive the wristband. The draw step demands this wristband
and won't proceed without it. So there's no way to draw against a notebook you
haven't just confirmed.

There's a nice extra safety catch built in. The wristband is \emph{single
use}---the moment the program draws, the wristband is torn off. You can't
save an old wristband from a previous update and reuse it. Each update must
earn a fresh proof. This stops a sneaky class of mistake where stale
permission gets reused after things have changed.

\textbf{Wristband two: ``the frozen lines still hold what we think they
hold.''} Remember the lines that have already scrolled off and frozen? The
program keeps notebook entries for some of them too. Normally it never touches
those---they're frozen, off-limits. But there's a subtle trap: what if the
program is about to \emph{switch} to a completely different conversation? Then
the new picture's frozen region should hold different text than what actually
froze earlier. If the program isn't careful, it can end up with the old frozen
lines stranded on screen while its notebook believes they've been replaced.

So there's a second check---and a second wristband---just for this: before
drawing a screen that has frozen lines above it, the program proves that those
frozen lines genuinely still match what it thinks. If they don't, it refuses to
draw normally and switches to repair instead.

\begin{pic}[Two doors, two guards]
Getting onto the screen is like passing through two doors. The first guard
checks that the \emph{part you can still see} matches your notebook. The
second guard checks that the \emph{part that already froze} matches too.
Between them, every line---visible or frozen---is vouched for before a single
mark is made.
\end{pic}

Two wristbands, two halves of the same guarantee: the part of the screen you
can still change, and the part you never can. Cover both, and there's no line
anywhere that could be wrong when it matters.

% ======================================================
\chapter{Delivering the message safely}
% ======================================================
\lettrine{O}{ne} more danger, and then our set of safeguards is complete. It's
about the actual moment of sending text to the screen, and it's sneakier than
it sounds.

Picture this. The program has figured out its changes and updated its
notebook to match the new screen. Then it starts sending the text to the
screen---and halfway through, the sending \emph{fails}. Maybe the connection
hiccupped (this program can even run over a remote link, like controlling a
computer across the internet). Now the screen only received \emph{half} the
changes, but the notebook already says the \emph{full} change happened. The
notebook and the screen have fallen out of step---the exact disaster
everything is built to prevent.

The fix is a rule about \emph{order} and \emph{atomicity}: the program treats
``send the text'' and ``update the notebook'' as a single, all-or-nothing
package. It hands over the new notebook \emph{only if} the send fully
succeeded. If the send fails partway, the half-baked new notebook is thrown
away entirely, and the program flips into a repair mode that assumes the
screen is now in an unknown state and rebuilds it cleanly.

\begin{pic}[Sign only on delivery]
Think of a courier delivering a parcel. You don't mark the parcel
``delivered'' the moment it leaves the van---you mark it delivered only when it's
actually in the customer's hands. If the van breaks down halfway, the parcel is
\emph{not} delivered, and you don't pretend it was. The program signs for the
new screen only once it's truly on the screen.
\end{pic}

\begin{insight}[The complete armour]
Put the pieces together. The program keeps a notebook shadowing the screen;
updates it in the same breath as every paint; double-checks the lines about to
freeze; verifies the notebook with a cheap fingerprint on every update; makes
those checks impossible to skip by requiring earned wristbands; and treats
``send'' and ``record'' as one all-or-nothing act. At no step can the notebook
quietly drift from the screen. The classic bugs aren't guarded against---they're
made impossible.
\end{insight}

% ======================================================
\chapter{When something still goes wrong}
% ======================================================
\lettrine{N}{o} matter how careful you are, some things are simply outside your
control. Another program might barge in and scribble on the screen. A remote
connection might drop and drown some text. The terminal window might get
resized mid-sentence. So the program can't only \emph{prevent} trouble---it
also needs a way to \emph{recover} when the world misbehaves.

It has three repair tools, and the whole trick is using the \emph{right} one.
They range from gentle to drastic:

\begin{description}
\item[The gentle one: quietly adjust the notebook.] Sometimes nothing is
  actually wrong on screen---the notebook just needs a small bookkeeping
  tweak to account for lines that have already frozen. This tool writes
  nothing to the screen at all. It just tidies the notebook. Cheapest, safest,
  used most often.
\item[The medium one: repaint what you can see.] This one redraws every line
  currently visible on screen, freshening it up. Good for clearing up leftover
  smudges inside the visible area. But it can only touch what's on screen---it
  can't reach the frozen lines above.
\item[The drastic one: wipe and start over.] This is the sledgehammer. It
  clears the whole screen \emph{and} the frozen history above, then repaints
  from a clean slate. It's the only tool that can reach the frozen lines---but
  it does so by destroying \emph{all} of them, including bits of your session
  from before the program even started. So it's reserved for moments when the
  user has clearly asked for a fresh start.
\end{description}

\begin{trap}[The tempting mistake]
Here's a trap the builders actually fell into. Suppose you switch to a
different conversation, and the old one had scrolled a lot of lines into the
frozen history. It's tempting to reach for the gentle or medium repair
tool---they feel safer. But neither can reach those old frozen lines, so a
ghostly copy of the previous conversation gets stranded above the new one. The
\emph{only} tool that actually cleans it up is the drastic wipe-and-restart.
Choosing the ``safe-feeling'' tool here is exactly the wrong call.
\end{trap}

The lesson is that matching the repair to the situation is its own skill. Use
the drastic wipe for a routine little update and you'll destroy the user's
history for no reason. Use a gentle tweak where a full wipe was needed and
you'll leave ghosts behind. The program has a small, careful chart matching
each kind of trouble to exactly the right tool---and a firm rule never to wire
the sledgehammer to anything that happens routinely.

% ======================================================
\chapter{Doing the sums exactly}
% ======================================================
\lettrine{O}{ne} last idea, and it's a lovely one because it's really a lesson
about honesty in engineering.

When the conversation gets very long, the program lets the oldest parts scroll
away into the frozen history and stops tracking them in detail---otherwise the
notebook would grow forever. But to let go of an old part cleanly, the program
has to know \emph{exactly} how many lines that old part took up on screen, so
its bookkeeping stays perfectly aligned. Off by even one, and a ghost line
appears.

Now, here's the mistake an earlier version made. The program tried to
\emph{guess} how tall each old part had been, by re-doing the arithmetic of how
the text would have wrapped. But that arithmetic was fiddly---it depended on
little margins and paddings---and every so often the guess was off by a line or
two. On narrow screens especially (imagine using this on a phone over a remote
connection), the guess drifted, and a duplicate ghost of the text appeared.

The fix is a principle worth carrying into life well beyond programming:

\begin{insight}[Don't guess what you can know]
Don't make one part of a system \emph{estimate} a number that another part
already \emph{knows} for certain. Have the part that actually did the work
write down the true number, and let everyone else simply read it.
\end{insight}

So instead of the bookkeeper guessing how tall each old block was, the part of
the program that \emph{actually drew} each block now records its true height on
the spot. Later, when it's time to let that block go, the bookkeeper just reads
the recorded number. No guessing, no drift, no ghost. The truth is measured
once, by the one who knows, and never reconstructed.

\begin{pic}[Measure once, at the source]
It's the difference between a builder guessing ``that wall was about nine feet''
from memory, versus reading the exact height he wrote on the blueprint the day
he built it. One drifts; the other is simply correct.
\end{pic}

% ======================================================
\chapter{The whole story, in one breath}
% ======================================================
\lettrine{L}{et's} stand back and see the whole thing at once, because now that
you have all the pieces, they form a single, satisfying shape.

The problem was: a program must keep a screen perfect, forever, even though it
can't look at the screen, can't fix lines once they scroll away, and lives in a
world where other things can barge in. That's a genuinely hard problem. The
solution came in four layers, each a different \emph{kind} of promise.

\begin{description}
\item[One notebook.] The program keeps exactly one private notebook shadowing
  the screen---never two, so nothing can disagree with itself.
\item[A faithful habit.] It updates that notebook in the same breath as every
  paint, and double-checks the lines about to freeze---so the notebook stays a
  true shadow, and the one irreversible moment is always guarded.
\item[A cheap constant check.] A little fingerprint lets it confirm the
  notebook is intact on every single update, without slowing down---so it can
  afford to always check.
\item[Unbreakable rules.] And it makes those checks impossible to skip, by
  requiring earned proof before any drawing or delivering happens---so the bug
  of ``forgetting to check'' cannot exist.
\end{description}

That progression is the real lesson, and it's more general than screens. It
goes from \emph{a tidy habit}, to \emph{a cheap way to catch yourself slipping},
to \emph{a design where slipping is impossible.} The people who built this kept
noticing the same thing: every past disaster was someone forgetting a check. So
they built a world where the check cannot be forgotten. That's the deepest idea
in the book.

And wrapped around all of it is a dose of humility: since the world can still
misbehave in ways you can't prevent, keep a small set of sharp repair tools,
and be disciplined about choosing the gentlest one that actually fixes the
problem.

\begin{insight}[What all this care buys]
You can chat with an AI in your terminal, watch a thousand-line answer stream
in, scroll up through your whole session, resize your window, switch
conversations---and never once see a duplicated box, a drifting line, or a
half-scrambled mess. Every frozen line is correct, because in the instant
before it froze, the program had \emph{proven}---not merely hoped---that it was
right.
\end{insight}

That's the entire story. A receipt printer you can't un-print. A magic notebook
you can never look away from to check. A last-second double-check at the
trapdoor. A fingerprint on the jar. Wristbands at the door. And a courier who
signs only on delivery. Put together, they let a little chat program do
something quietly remarkable: draw on a shared, unforgiving screen, thousands of
times a second, and never, ever leave a scar.

\vspace{1cm}
\begin{center}
{\sffamily\color{rulegray}\rule{0.3\linewidth}{0.5pt}}\\[6pt]
{\itshape\color{accent} --- the end ---}
\end{center}

% ---------------- BACK MATTER ----------------
\backmatter
\chapter*{A little glossary}
\addcontentsline{toc}{chapter}{A little glossary}
\markboth{A little glossary}{}

In case you want the everyday nicknames matched to the real names
programmers use:

\begin{description}
\item[The receipt roll / frozen history] The terminal's \emph{scrollback}---the
  saved lines you reach by scrolling up. Once a line enters it, the program
  can't change it.
\item[Drawing in line, not taking over] \emph{Inline} rendering, as opposed to
  a \emph{full-screen} program that owns the whole display.
\item[The notebook / shadow] The renderer's private copy of what it believes is
  on screen. It's compared against the freshly-laid-out new picture to decide
  what to redraw.
\item[The compare] The \emph{diff}---finding where the notebook and the new
  picture differ, so only the changes get drawn.
\item[The pen] The \emph{cursor}: the invisible position where text appears.
  The program tracks it by memory, never by looking.
\item[The last-second double-check] Forcing a full repaint of any line about to
  scroll off, so a stale line can never freeze.
\item[The fingerprint] A \emph{hash}: one small number summarising the notebook,
  so corruption is caught instantly and cheaply.
\item[The wristbands] Proof tokens the code must hold before it's allowed to
  draw or deliver---making the safety checks impossible to skip.
\item[Sign only on delivery] Updating the notebook only after the screen write
  fully succeeds, so a failed send can't desync the two.
\item[The three repair tools] From gentle to drastic: adjust-the-notebook,
  repaint-the-visible, and wipe-and-restart.
\end{description}

\vspace{6pt}
\noindent This book explains, in plain language, the real design used by a
terminal AI assistant to keep its screen perfectly in sync---the same ideas the
programmers rely on, just told with receipts and notebooks instead of code.

\end{document}
